\documentclass[10pt,a4paper]{amsart}
\usepackage[margin=19mm]{geometry}
\usepackage{amsmath,amssymb,mathtools}
\usepackage[scr=rsfs,cal=euler]{mathalfa}
\usepackage{xfrac}
\usepackage[unicode,hidelinks,bookmarks=false]{hyperref}
\newcommand{\ie}{i.e.\ }
\newcommand{\cf}{cf.\ }
\newcommand{\resp}{respectively\ }
\newcommand{\Subsection}{\subsection}
\providecommand{\nobreakdash}{\nobreak\hbox{-}}
\setlength{\emergencystretch}{2em}
\multlinegap0pt
\usepackage{float}
\usepackage{mathtools}
\usepackage{amssymb}
\usepackage{bbm}
\usepackage[shortlabels]{enumitem}
\usepackage[normalem]{ulem}
\usepackage{natbib}
\usepackage{xcolor}
\newtheorem{theorem}{Theorem}
\newtheorem{proposition}{Proposition}
\newtheorem{corollary}{Corollary}
\newcommand{\CProb}[2]{\pr\left(#1 \Bigm| #2\right)}	%	Conditional probability command, 
\newcommand{\E}{\mathbb{E}}								%	short hand for blackboard E
\newcommand{\Exp}[1]{\E\left[#1\right]}					%	Standard expectation command, argument
\newcommand{\Prob}[1]{\pr\left(#1\right)}					%	Standard probability command, 
\newcommand{\R}{\mathbb R}
\newcommand{\T}{\mathbb T}
\newcommand{\bT}{\mathbf T}
\newcommand{\cB}{\mathcal B}
\newcommand{\eps}{\epsilon}
\newcommand{\floor}[1]{\left\lfloor #1\right\rfloor}
\newcommand{\pr}{\mathbb{P}}								%	short hand for blackboard P
\newcommand{\toinp}{\overset{\mathbb P}{\longrightarrow}}
\title{On exponentially height-penalized random trees}
\author{Louigi Addario-Berry, Beno\^it Corsini, Neeladri Maitra, Meltem \"Unel}
\date{Source: arXiv:2512.17747v1 (CC BY 4.0)}
\begin{document}
\maketitle

\noindent\emph{Source and licence.} On exponentially height-penalized random trees. Version arXiv:2512.17747v1 (CC BY 4.0), distributed by arXiv under the Creative Commons Attribution 4.0 International licence, http://creativecommons.org/licenses/by/4.0/, as recorded in the arXiv licence metadata for this exact version and shown on the abstract page https://arxiv.org/abs/2512.17747v1. Full licence notice: \texttt{licenses/arxiv-2512.17747-CC-BY-4.0.txt}.\par\smallskip

\noindent\textit{Editorial scope.} Counted unit: the article's theorem thm:width (source0.tex lines 505--512). The article's complete original proof of it is delivered with it (source0.tex lines 1610--1683, one \texttt{proof} environment of 74 lines, which the article places away from the statement). No mathematics is written, repaired or re-derived: the statement and the proof are the article's own lines, in the article's own notation, and no retained line is substituted or re-flowed.  Retained uncounted as the source-local context the proof needs: lines 434--441; lines 1469--1475; lines 1525--1531; lines 1567--1573; lines 1603--1608.  One declared adaptation: exactly 1 ancillary strings inside the retained spans are omitted (image-inclusion commands and, where the source repeats a label, the repeated label definitions) and no other line differs from the source; the figure environments, with their placement, captions and labels, are kept, and the package records the omitted strings in fidelity receipts bound to the affected bodies.  No retained line contains a citation, so the wrapper carries no bibliography and no bibliography entry.  The retained lines contain the article's own diagram code, which the wrapper typesets with the standard tikz packages; no image file is delivered and the package declares no asset dependency.  Editorial formatting only, in addition: an AMS \texttt{amsart} preamble that restates the article's own declarations and macros used by the retained lines, namely the environments \texttt{theorem}, \texttt{proposition}, \texttt{corollary} on separate counters, together with the standard packages those lines need.  The retained environments print in this document as Theorem 1, Proposition 1, Corollary 1 in order of appearance, whereas the article prints the same items with the numbering of its own full document.  The complete native source of the article as distributed in the arXiv e-print for this exact version is bundled unchanged as \texttt{sources/191342/source0.tex}.
\begin{theorem}[Asymptotic height in the intermediate regime]\label{thm:inter lln}
    Let $\mu=\mu_n$ be such that $1/\sqrt{n}\ll\mu\ll n$.
    Let $\mathbf{T}_n$ be a $\mu$-height-biased tree of size $n$.
    Then
    \begin{align*}
        \frac{h(\mathbf{T}_n)}{\left(\frac{2\pi^2n}{\mu}\right)^{1/3}}\toinp 1\,.
    \end{align*}
\end{theorem}

\begin{proposition}\label{prop:contour_random_narrow_exc_pos}
    Let $\mu=\mu_n$ such that $1/\sqrt{n}\ll\mu\ll n$.
    Then for any $\epsilon>0$, there exists $\delta>0$ such that, for any sequence of events $(A_n)_{n\geq1}$, we have
    \begin{align*}
        \limsup_{n\rightarrow\infty}\frac{1}{(\mu^2n)^{1/3}}\log\left(\frac{\Prob{\bT_n\in A_n,\left|\frac{h(\bT_n)}{(2\pi^2n/\mu)^{1/3}}-1\right|\leq\delta}}{\CProb{\T_n\in A_n}{\left|\frac{h(\T_n)}{(2\pi^2n/\mu)^{1/3}}-1\right|\leq\delta}}\right)\leq\epsilon\,.
    \end{align*}
\end{proposition}

\begin{proposition}[Subgaussian width tails for bridges] \label{prop:bridge_subg_tail}
    Let $B^{(n)}_x=(B^{(n)}_{x;t}:0\leq t \leq n)$ be a uniform bridge of length $n$ ending at $x$, that is a uniform element from $\cB^{(n)}_x$.
    Then, for any $A>0$, there exists $R_1,R_2>0$ such that, for any $|x|\leq A\sqrt{n}$ and $T\geq0$, we have
    \begin{align*}
        \Prob{\frac{w\big(B^{(n)}_x\big)}{\sqrt{n}}>T}\leq R_1e^{-R_2T^2}\,.
    \end{align*}
\end{proposition}

\begin{corollary}[Exponential moments of the width of uniform bridges]\label{cor:exp_mom_local_time}
    We denote by $B^{(n)}_{x,y}$ a uniform bridge of length $n$ from $x$ to $y$, that is a uniform element from $\cB^{(n)}_{x,y}$.
    Then, for any $s\in\R$ and $A>0$, there exists $M>0$ such that, for all $n\geq1$ and all $0\leq x,y\leq A\sqrt{n}$, we have
    \begin{align*}
        \Exp{\exp\left(s\frac{w\big(B^{(n)}_{x,y}\big)}{\sqrt{n}}\right)}\leq M\,.
    \end{align*}
\end{corollary}

\begin{figure}[H]
    \centering
    
    \caption{A representation of the multiscale decomposition used for the proof of Theorem~\ref{thm:width}. We bound the height with $h_n^\delta=(1+\delta)(2\pi^2n/\mu)^{1/3}$ and split the contour of a uniform tree into $r\sim(\mu^2n)^{1/3}$ bridges of length $\Delta_t\sim2(n/\mu)^{2/3}$. This way the length of each bridge is of the same order as the square of its fluctuations, thus are objects from the \emph{Brownian world}, admitting good enough width tails for our purposes, thanks to Proposition~\ref{prop:bridge_subg_tail}.}
    \label{fig:broscale}
\end{figure}

\begin{theorem}[Width in the non-Brownian regime]\label{thm:width}
Let $\mu=\mu_n$ be such that $\mu\gg (\log n)^{3/2}/\sqrt{n}$.
Let $\mathbf{T}_n$ be a $\mu$-height-biased tree of size $n$.
Then for any $\eps>0$, we can find $K=K(\eps)>0$ large enough, such that
        \begin{align*}
            \liminf_{n \to \infty} \Prob{K^{-1}\leq\frac{w(\mathbf{T}_n)}{n\wedge\left((\mu n^2)^{1/3}\right)}\leq K}\geq 1-\eps\,.
        \end{align*}
\end{theorem}

\begin{proof}[Proof of Theorem~\ref{thm:width}]
    To start the proof, first note that, since the width counts the number of nodes at a certain height in the tree, by summing over all possible heights we directly have the bound
    \begin{align*}
        w(\bT_n)\cdot h(\bT_n)\geq n\,.
    \end{align*}
    Using the result of Theorem~\ref{thm:inter lln} which states that $h(\bT_n)$ is asymptotically of order $(n/\mu)^{1/3}$, we directly see that $w(\bT_n)$ is at least of order $(\mu n^2)^{1/3}$ as desired.
    This proves the lower bound of the theorem and we now focus on proving that $w(\bT_n)$ is not of higher order than $(\mu n^2)^{1/3}$.
    Following the proof of Proposition~\ref{prop:contour_random_narrow_exc_pos}, we let $h_n=(2\pi^2n/\mu)^{1/3}$ be the asymptotic height of $\bT_n$.

    Fix for now some $K>0$ and $\epsilon>0$.
    We start by applying Proposition~\ref{prop:contour_random_narrow_exc_pos} to the event $A_n=\{w(T)>K(\mu n^2)^{}\}$ to find $\delta>0$ such that
    \begin{align*}
        \limsup_{n\rightarrow\infty}\frac{1}{(\mu^2n)^{1/3}}\log\left(\frac{\Prob{w(\bT_n)>K(\mu n^2)^{1/3},\left|\frac{h(\bT_n)}{h_n}-1\right|\leq\delta}}{\CProb{w(\T_n)>K(\mu n^2)^{1/3}}{\left|\frac{h(\T_n)}{h_n}-1\right|\leq\delta}}\right)\leq\epsilon\,.
    \end{align*}
    By further using Theorem~\ref{thm:inter lln}, we know that $h(\bT_n)$ converges in probability to $h_n$.
    Combining this with the previous result leads to
    \begin{align}\label{eq:bound width}
        \Prob{\Big.w(\bT_n)>K(\mu n^2)^{1/3}}\leq e^{O(\epsilon(\mu^2n)^{1/3})}\CProb{w(\T_n)>K(\mu n^2)^{1/3}}{\left|\frac{h(\T_n)}{h_n}-1\right|\leq\delta}+o(1)\,.
    \end{align}
    Proving the theorem now boils down to proving an upper bound on the width of a uniform tree.

    Given a uniform tree $\T_n$, recall we let $C^{(\T_n)}=(C^{(\T_n)}_t:0\leq t\leq 2n-1)$ to be its contour walk.
    We are going to split this walk into a sequence of bridges of length $2(n/\mu)^{2/3}$ so that, once conditioned to be within the range $[0,h_n)$, their length is exactly this range squared. 
    We refer to Figure~\ref{fig:broscale} for a representation of this idea.

    Let $r=r_n=\floor{(\mu^2n)^{1/3}}$ and set $t_0,\ldots,t_r$ to evenly spaced time points from the interval $2n-1$, meaning that
    \begin{align*}
        t_i=\floor{\frac{i(2n-1)}{r}}\,.
    \end{align*}
    We observe here that $r\sim(\mu^2n)^{1/3}$ and that $\Delta_t=\max\{t_i-t_{i-1}\}\sim2(n/\mu)^{2/3}$.
    Finally, we denote by $B_i=(C^{(\T_n)}_{t+t_{i-1}}:0\leq t\leq t_i-t_{i-1})$ the walk corresponding to the portion of $C^{(n)}$ between time $t_{i-1}$ and time $t_i$ (rescaled in time to start at $0$).
    We now make two important observations.
    \begin{enumerate}
        \item Since $C^{(\T_n)}$ is obtained by concatenating $B_1,\ldots,B_r$, its width is bounded by the sum of their widths:
        \begin{align*}
            w\big(C^{(\T_n)}\big)\leq\sum_{i=1}^rw(B_i)\,.
        \end{align*}
        \item By the Markov property $B_1,B_2,\ldots,B_r$ are independent and uniform bridges conditionally given $C^{(\T_n)}_{t_i}$ for $0\leq i\leq r$.
    \end{enumerate}
    To use these observations, start by applying a standard exponential Markov's inequality to see that, for any $\rho>0$, we have
    \begin{align*}
        \CProb{w(\T_n)>K(\mu n^2)^{1/3}}{\left|\frac{h(\T_n)}{h_n}-1\right|\leq\delta}
        &\leq e^{-\rho K(\mu n^2)^{1/3}}\Exp{e^{\rho w(\T_n)}~\bigg|~\left|\frac{h(\T_n)}{h_n}-1\right|\leq\delta}\,.
    \end{align*}
    We then combine the first observation and a nested expectation to obtain that
    \begin{align*}
        &\hspace{-0.5cm}\CProb{w(\T_n)>K(\mu n^2)^{1/3}}{\left|\frac{h(\T_n)}{h_n}-1\right|\leq\delta}\\
        &\leq e^{-\rho K(\mu n^2)^{1/3}}\Exp{\Exp{e^{\rho w(\T_n)}~\bigg|~\left|\frac{h(\T_n)}{h_n}-1\right|\leq\delta,C^{(\T_n)}_{t_0},C^{(\T_n)}_{t_1},\ldots,C^{(\T_n)}_{t_r}}}\\
        &\leq e^{-\rho K(\mu n^2)^{1/3}}\Exp{\Exp{e^{\rho\left[w(B_1)+\ldots+w(B_r)\right]}~\bigg|~\left|\frac{h(\T_n)}{h_n}-1\right|\leq\delta,C^{(\T_n)}_{t_0},C^{(\T_n)}_{t_1},\ldots,C^{(\T_n)}_{t_r}}}\,.
    \end{align*}
    The condition that $h(\T_n)$ is bounded can be reduced to simply having $0\leq C^{\T_n}_{t_i}\leq(1+\delta)h_n$ by dividing with the probability that a walk of length $2n-1$ is a contour whose height is of order $h_n$, which itself is of order $e^{O(\log n)+O((\mu^2n)^{1/3})}$.
    Finally, the second observation states that, under the above conditioning, all bridges are independent from each other and their endpoints both belong to the interval $[0,(1+\delta)h_n]$.
    Further using that the length of the bridges is bounded by $\Delta_t\sim2(n/\mu)^{2/3}$ and that extending a walk can only increase its width, the previous bound can be transformed into
    \begin{align*}
        \CProb{w(\T_n)>K(\mu n^2)^{1/3}}{\left|\frac{h(\T_n)}{h_n}-1\right|\leq\delta}
        &\leq e^{-\rho K(\mu n^2)^{1/3}+O(\log n)+O((\mu^2n)^{1/3})}\left(\max_{0\leq x,y\leq A\sqrt{\Delta_t}}\Exp{e^{\rho w\left(B^{(\Delta_t)}_{x,y}\right)}}\right)^r\,,
    \end{align*}
    where $A>0$ is some constant such that $(1+\delta)h_n\leq A\sqrt{\Delta_t}$ and $B^{(\Delta_t)}_{x,y}$ is a uniform random bridge of length $\Delta_t$ going from $x$ to $y$.

    For the final part of this proof, start by replacing $\rho$ with $1/\sqrt{\Delta_t}$ and apply Corollary~\ref{cor:exp_mom_local_time} with the same $A$ and $s=1$ to find some $M>0$ such that, for any $\Delta_t>0$
    \begin{align*}
        \max_{0\leq x,y\leq A\sqrt{\Delta_t}}\Exp{e^{w\left(B^{(\Delta_t)}_{x,y}\right)/\sqrt{\Delta_t}}}\leq M\,.
    \end{align*}
    Finally, by recalling that $\mu\gg(\log n)^{3/2}/\sqrt{n}$, that $r\sim(\mu^2n)^{1/3}$, and that $\Delta_t\sim2(n/\mu)^{2/3}$, combining the previous two bounds lead to
    \begin{align*}
        \CProb{w(\T_n)>K(\mu n^2)^{1/3}}{\left|\frac{h(\T_n)}{h_n}-1\right|\leq\delta}&\leq\exp\left(-\big(1+o(1)\big)\left[\frac{K}{\sqrt{2}}-\log M+O(1)\right](\mu^2n)^{1/3}\right)\,.
    \end{align*}
    Plugging this back into~\eqref{eq:bound width} finally leads to
    \begin{align*}
        \Prob{\Big.w(\bT_n)>K(\mu n^2)^{1/3}}\leq\exp\left(-\big(1+o(1)\big)\left[\frac{K}{\sqrt{2}}-\log M+O(\epsilon)+O(1)\right](\mu^2n)^{1/3}\right)+o(1)
    \end{align*}
    Observe that the definitions of $M$, $O(\epsilon)$, and $O(1)$ only depend on $\epsilon$, $\delta$, and $A$ and are thus independent of $K$.
    This means that we can take $K$ large enough such that the exponential converges to $0$, thus proving the desired upper bound and concluding the proof.
\end{proof}

\end{document}
